\section{Introduction}
\label{sec:introduction}

Large Language Models (LLMs) have become the default substrate for
modern natural language processing, but the prevailing usage pattern
treats each fine-tuned model as a single static artifact. In
practice, downstream tasks rarely arrive simultaneously: a model
deployed in a research group, startup, or product team is repeatedly
asked to acquire \emph{new} capabilities (a new domain, a new
language, a new instruction style) without losing the capabilities it
already possesses. This is the classical setting of
\emph{continual learning}, and the failure mode it exposes---a sharp
collapse in performance on previously mastered tasks the moment a new
task is introduced---is \emph{catastrophic forgetting} (CF)
\cite{mccloskey1989catastrophic,luo2023empirical}.

The recent CF-mitigation literature for LLMs is wide and
fast-moving. Regularization-based methods penalize movement of
important parameters \cite{kirkpatrick2017overcoming,li2024revisiting,song2025hierarchical};
rehearsal-based methods reintroduce data
from earlier tasks during current-task training \cite{huang2024ssr,reynolds2025mixed};
architectural and PEFT-based methods isolate
task-specific knowledge in dedicated low-rank subspaces
\cite{hu2021lora,wang2023olora,fawi2024curlora,nayak2025sculpting}.
Despite the breadth of techniques, the experimental setting in which
they are typically validated is remarkably homogeneous: a
$\geq$7B-parameter base model, an A100-class GPU, and large
fine-tuning budgets. The resulting body of evidence is largely
silent on the regime that most students, independent researchers,
and small teams actually inhabit: a sub-2B-parameter base model and a
single free-tier GPU (e.g.\ an NVIDIA T4 on Colab or Kaggle).

This work fills that gap. We perform a controlled empirical
comparison of six representative CF-mitigation strategies on small,
publicly available instruction-tuned LLMs under a strict
free-tier compute envelope. Concretely, we make the following
contributions:

\begin{itemize}
    \item We establish a reproducible \emph{small-model,
    free-tier} benchmark for CF on a four-task sequence drawn from
    GLUE (SST-2 $\rightarrow$ MRPC $\rightarrow$ CoLA $\rightarrow$
    MNLI), with all six methods implemented uniformly on top of
    QLoRA 4-bit fine-tuning.

    \item We formalize \emph{Minimal Mixed Rehearsal} (MMR), a simple
    rehearsal recipe that retains only a small fraction
    ($\sim$\todo{$f\%$}) of each earlier task's training set and
    interleaves it with the current task at every optimization step.
    The recipe is parameter-free at inference time and adds
    negligible memory and wall-clock overhead.

    \item We report not only the standard continual-learning metrics
    (Average Accuracy, Forgetting, Backward Transfer) but also
    \emph{wall-clock training time} and \emph{peak GPU memory}, and
    visualize the resulting accuracy/compute trade-off as a Pareto
    frontier. This compute-aware view is, to our knowledge, missing
    from existing comparative CF studies.

    \item Across both Qwen2.5-1.5B-Instruct and Phi-3.5-mini-instruct,
    we observe that
    \todo{state the headline empirical finding once results are in,
    e.g.\ ``MMR matches O-LoRA in average accuracy while reducing
    training time by X\%''}.

    \item We release code, configurations, and run logs to enable
    end-to-end reproduction of every reported number on a single
    Colab T4 session.
\end{itemize}

The remainder of the paper is organized as follows.
Section~\ref{sec:related} situates our work in the existing CF
literature. Section~\ref{sec:method} introduces our notation,
formalizes each comparison method, and presents MMR.
Section~\ref{sec:setup} details the experimental protocol.
Section~\ref{sec:results} reports our main empirical results.
Section~\ref{sec:discussion} draws practical guidance from the
findings, and Section~\ref{sec:conclusion} concludes.
